\subsection{Loop-disentanglable spaces}

DEFINITION 2. --- A topological space B is said to be loop-disentanglable if it is locally path-connected and if every point b of B has a neighbourhood V such that the homomorphism from \(\pi_1(V,b)\) into \(\pi_1(B,b)\) deduced from the canonical injection has as its image the identity element.

A locally path-connected topological space is loop-disentanglable if and only if each of its connected components is.

Let B be a locally path-connected topological space. Suppose that every point of B has a neighbourhood V which is loop-disentanglable. Then B is loop-disentanglable.

Every locally path-connected and simply path-connected topological space is loop-disentanglable. This is in particular the case for a locally path-connected space homotopic to a point.

Remarks. --- 1) There exist topological spaces B, connected and locally path-connected, such that certain points a of B, but not all, have a neighbourhood V such that the homomorphism from \(\pi_{1}(V,a)\) into \(\pi_{1}(B,a)\) is trivial (III, p. 336, exerc. 6).

\begin{enumerate}
\def\labelenumi{\arabic{enumi})}
\setcounter{enumi}{1}
\tightlist
\item
  Let B be a loop-disentanglable space; every operation of the groupoid \(\varpi(\mathrm{B})\) on a B-space is without local monodromy (cf. III, p. 313, remark). In particular, for a map \(p: E \rightarrow B\) to make E a covering of B, it is necessary and sufficient that it be étale, separated, and that it satisfy the path-lifting property (III, p. 315, corollary 3).
\end{enumerate}

PROPOSITION 1. --- The product space of a finite family of loop-disentanglable spaces is loop-disentanglable.

It suffices to prove that, if A and B are two loop-disentanglable spaces, then their product A × B is also. Under these conditions, the space A × B is indeed locally path-connected (III, p. 261, prop. 9). Let (a, b) be a point of A × B. By hypothesis there exists a neighbourhood U of a (resp. a neighbourhood V of b) such that the image of the homomorphism \(i_* \colon \pi_1(\mathrm{U}, a) \to \pi_1(\mathrm{A}, a)\) deduced from the canonical injection \(i \colon \mathrm{U} \to \mathrm{A}\) (resp. of the homomorphism \(j_* \colon \pi_1(\mathrm{V}, b) \to \pi_1(\mathrm{B}, b)\) deduced from the canonical injection \(j \colon \mathrm{V} \to \mathrm{B}\) ) is reduced to the identity element. Then U × V is a neighbourhood of (a, b) in A × B. The homomorphism \(\pi_1(\mathrm{U} \times \mathrm{V}, (a, b)) \to \pi_1(\mathrm{A} \times \mathrm{B}, (a, b))\) is identified with the homomorphism \((i_*, j_*)\) (III, p. 297, corollary of prop. 4). Its image is therefore reduced to the identity element. This proves that A × B is loop-disentanglable.

PROPOSITION 2. --- a) Every covering of a loop-disentanglable space is loop-disentanglable.

\begin{enumerate}
\def\labelenumi{\alph{enumi})}
\setcounter{enumi}{1}
\item
  Conversely, let \(p \colon \mathrm{E} \to \mathrm{B}\) be an étale and surjective map and suppose that \(\mathrm{E}\) is a loop-disentanglable space. Then \(\mathrm{B}\) is loop-disentanglable.
\item
  Let B be a loop-disentanglable topological space and let \((\mathrm{E}, p)\) be a covering of B. The space E is then locally path-connected. Let x be a point of E, write \(b = p(x)\) , and let V be a neighbourhood of b in B such that the canonical homomorphism \(\pi_{1}(\mathrm{V}, b) \to \pi_{1}(\mathrm{B}, b)\) is trivial. Write \(i: V \to B\) and \(j: p^{-1}(V) \to E\) for the canonical injections. By hypothesis, the image of the homomorphism \(\pi_{1}(i, b)\) is reduced to the identity element. Since \(p \circ i = j \circ p\) and \(\pi_{1}(p, x)\) is injective, the image of the homomorphism \(\pi_{1}(j, x)\) is reduced to the identity element. This proves that E is loop-disentanglable.
\item
  The space B is locally path-connected. Let b be a point of B and let x be a point of \(E_{b}\) . Since p is étale, there exists a neighbourhood U of x in E such that p induces a homeomorphism of U onto \(p(\mathrm{U})\) . Let V be a neighbourhood of x in E contained in U such that the homomorphism \(\pi_{1}(j,x)\) is trivial, where \(j:V\to E\) denotes the canonical injection. Write \(q:V\to p(V)\) for the map deduced from p by passing to subspaces; it is a homeomorphism. Also write i for the canonical injection of \(p(\mathrm{V})\) into B. We have \(p\circ j=i\circ q\) , therefore \(\pi_{1}(i,b)\circ\pi_{1}(q,x)=\pi_{1}(p,x)\circ\pi_{1}(j,x)\) is the trivial homomorphism. Since the homomorphism \(\pi_{1}(q,x)\) is an isomorphism, the homomorphism \(\pi_{1}(i,b)\) is trivial. It follows that the space B is loop-disentanglable.
\end{enumerate}

PROPOSITION 3. --- Let B be a loop-disentanglable topological space. Let (Y, q) be a covering of B and let (Z, p) be a covering of Y. The topological space Z, equipped with the map \(q \circ p\) , is then a covering of B.

Indeed, the map \(q \circ p\) is étale (I, p. 29, prop. 6) and separated (I, p. 25, prop. 2). By remark 2 above, it therefore suffices to show that it satisfies the path-lifting property. Let \(z\) be a point of Z and let \(c \colon \mathbf{I} \to \mathbf{B}\) be a path in B with origin \(q(p(z))\) . There exists a path \(c'\) with source \(p(z)\) in Y which lifts \(c\) , since Y is a covering of B (III, p. 302, cor. 2). Since Z is a covering of Y, there exists a path \(c''\) with source \(z\) in Z which lifts \(c'\) ; the path \(c''\) lifts \(c\) . This proves the proposition.

\subsection{Universal covering of a loop-disentanglable space}

Let B be a topological space and let b be a point of B. Recall that \(\Lambda_{b}(\mathrm{B})\) denotes the subspace of \(\mathcal{C}_{c}(\mathbf{I};\mathrm{B})\) formed by the paths with origin b, equipped with the quotient topology of the compact-convergence topology. We denote by \(e_{\mathrm{B}}\colon\Lambda_{b}(\mathrm{B})\to\mathrm{B}\) the map which, to a path, associates its endpoint; it is continuous. Write \(\lambda_{b}(\mathrm{B})\) for the quotient space of \(\Lambda_{b}(\mathrm{B})\) by the strict homotopy relation and equip it with the quotient topology. Since two strictly homotopic paths have the same endpoint, the map \(e_{B}\) defines, by passing to the quotient, a continuous map \(\varepsilon_{\mathrm{B}}\colon\lambda_{b}(\mathrm{B})\to\mathrm{B}\) , called the endpoint map.

THEOREM 1. --- Let B be a connected and locally path-connected topological space and let b be a point of B. The following properties are equivalent:

\begin{enumerate}
\def\labelenumi{(\roman{enumi})}
\item
  The space B is loop-disentanglable.
\item
  There exists a non-empty covering of B which is simply path-connected.
\item
  The space \(\lambda_{b}(\mathrm{B})\) , equipped with the endpoint map \(\varepsilon_{\mathrm{B}}\colon\lambda_{b}(\mathrm{B})\to\mathrm{B}\) , is a covering of B.
\item
  The group \(\pi_1(B, b)\) is discrete for the admissible topology.
\item
  The group \(\pi_{1}(B,b)\) is discrete for the quotient topology of the compact convergence topology.
\end{enumerate}

Moreover, when these conditions are satisfied, the space \(\lambda_{b}(\mathrm{B})\) is simply path-connected, simply connected, Galois with group \(\pi_{1}(\mathrm{B},b)^{\circ}\) , and the pointed covering \((\lambda_{b}(\mathrm{B}),\varepsilon_{b})\) is a universal covering of the pointed space \((\mathrm{B},b)\) .

It follows from the definition of a loop-disentanglable space that the subgroup of \(\pi_{1}(B,b)\) reduced to the identity element is admissible if and only if B is loop-disentanglable. This proves the equivalence (i)⇔(iv). Moreover, the equivalence of properties (iv) and (v) follows from remarks 4 and 5 of III, p. 320.

(iv)⇒(ii). By III, p. 316, prop. 5, there exists a covering (E, p) of B and a point \(x \in E_{b}\) such that \(p_{*}(\pi_{1}(E, x)) = \{e\}\) ; in particular, \(\pi_{1}(E, x)\) is reduced to the identity element. The arc-connected component of x in E is then a nonempty covering of B (I, p. 80, cor. 1 of prop. 6), simply path-connected.

(ii)⇒(iii). Let E be a nonempty simply path-connected covering of B. Let x be a point of \(E_{b}\) (such exist because B is connected, I, p. 74, prop. 4). The projection \(q: E \to B\) induces a homeomorphism \(\Lambda_{x}(E) \to \Lambda_{b}(B)\) (III, p. 302, cor. 2 of prop. 3), whence, by passing to arc-connected components, a homeomorphism \(q_{*}: \lambda_{x}(E) \simeq \lambda_{b}(B)\) . The endpoint map \(\Lambda_{x}(E) \to E\) is continuous and open, since E is locally path-connected (III, p. 264, corollary). The map \(\varepsilon_{\mathrm{E}}: \lambda_{x}(E) \to E\) that it defines by passing to arc-connected components is therefore continuous and open. The map \(\varepsilon_{\mathrm{E}} \circ (q_{*})^{-1}: \lambda_{b}(B) \to E\) is bijective, continuous, and open. This proves that the topological space \(\lambda_{b}(B)\) , endowed with the topology of compact convergence and the endpoint map, is a covering of B.

(iii)⇒(v). Indeed, the set \(\pi_{1}(B,b)\) , endowed with the quotient topology of the topology of compact convergence, is identified with the fiber of the B-space \(\lambda_{b}(B)\) over b, at the class of the constant loop at b. If \(\lambda_{b}(B)\) is a covering of B, then \(\pi_{1}(B,b)\) is discrete for the topology of compact convergence.

Suppose these assertions hold. By the foregoing, the space \(\lambda_b(B)\) is then a covering of B that is simply path-connected, hence simply connected. The pointed space ( \(\lambda_b(B), \varepsilon_b\) ) is consequently a universal covering of (B, \(b\) ) (I, p. 126, corollary of prop. 3) and the covering \(\lambda_b(B)\) is a Galois covering of B (I, p. 120, prop. 1). The canonical homomorphism \(h_{(\lambda_b(B), \varepsilon_b)}: \pi_1(B, b) \to \text{Aut}_B(\lambda_b(B))\) is then an isomorphism (III, p. 306, prop. 5).

Remark 1. --- Let B be a connected loop-disentanglable space and let b be a point of B. It follows from Theorem 1, (iv), that every action of the group \(\pi_{1}(B,b)\) is admissible. Consequently, every \(\pi_{1}(B,b)\) -set is isomorphic to a \(\pi_{1}(B,b)\) -set \(E_{b}\) , where E is a covering of B (III, p. 318, prop. 7). Recall also (III, p. 310, prop. 2) that if E and \(E'\) are coverings of B, the map \(f \mapsto f_{b}\) induces a bijection from the set of morphisms of B-spaces from E into \(E'\) onto the set of morphisms of \(\pi_{1}(B,b)\) -sets from \(E_{b}\) onto \(E'_{b}\) .

In the language of categories, one says that the functor which, to a covering E of B, associates the fiber \(E_b\) is an equivalence of categories from the category of coverings of B to the category of \(\pi_1(B,b)\) -sets.*

COROLLARY 1. --- Let B be a loop-disentanglable topological space and let b be a point of B. Let E be a connected covering of B and let x be a point of the fiber \(E_{b}\) . The following properties are equivalent:

\begin{enumerate}
\def\labelenumi{(\roman{enumi})}
\item
  The group \(\pi_{1}(\mathrm{E},x)\) is trivial.
\item
  The space E is simply connected.
\item
  The pointed space (E, x) is a universal covering of (B, b).
\end{enumerate}

The implication (i)⇒(ii) has already been proved under the sole hypothesis that the space B is connected and locally path-connected (III, p. 309, cor. 2 of prop. 1), and the implication (ii)⇒(iii) without any hypothesis (I, p. 126, corollary).

Let us prove (iii)⇒(i). By Theorem 1 of IV, p. 342, there exist a covering E' of B and a point \(x'\) of the fiber \(E'_{b}\) such that the group \(\pi_1(E', x')\) is reduced to the identity element. Under hypothesis (iii), there exists

a pointed B-morphism \(f \colon (\mathrm{E}, x) \to (\mathrm{E}', x')\) . Denote by \(p \colon \mathrm{E} \to \mathrm{B}\) and \(p' \colon \mathrm{E}' \to \mathrm{B}\) the projections of coverings \(\mathrm{E}\) and \(\mathrm{E}'\) ; we have \(p = p' \circ f\) therefore \(\pi_1(p, x) = \pi_1(p', x') \circ \pi_1(f, x)\) . Since the group \(\pi_1(\mathrm{E}', x')\) is trivial, the injective homomorphism \(\pi_1(p, x)\) has as its image the identity element. This proves that the group \(\pi_1(\mathrm{E}, x)\) is reduced to the identity element.

COROLLARY 2. --- Let B be a loop-disentanglable topological space and let b be a point of B. Let (E, x) be a universal covering of the pointed space (B, b). For every covering E' of B and every point \(x' \in E'_{b}\), the unique B-morphism f: E → E' such that f(x) = x' makes E a covering of E'. The pointed space (E, x) is then a universal covering of (E', x').

By prop. 7 of I, p. 81, the space E, equipped with the map f, is a covering of E'. The last assertion then follows from corollary 1.

Let B be a loop-disentanglable topological space and let \(a\) , \(b\) be two points of B. The quotient topological space \(\varpi_{a,b}(\mathrm{B})\) is homeomorphic to the space \(\pi_1(\mathrm{B},a)\) , equipped with the quotient topology of the compact-convergence topology (III, p. 293, remark 3), hence is discrete by theorem 1 of IV, p. 342. Consequently, the strict homotopy class of every path in B joining \(a\) to \(b\) is an open subset of \(\Lambda_{a,b}(\mathrm{B})\) . More generally:

PROPOSITION 4. --- Let B be a loop-disentanglable topological space. Strict homotopy is an open equivalence relation in the topological space \(\mathcal{C}_{c}(\mathbf{I};\mathbf{B})\) .

Suppose first that the space B is simply path-connected. Let \(\varphi\colon\Lambda(B)\to B\times B\) be the map which, to a path \(c\) in B, associates the pair \((c(0),c(1))\) consisting of its origin and its endpoint. For two paths in B to be strictly homotopic, it is necessary and sufficient that they have the same origin and the same endpoint (III, p. 292, prop. 2). Since the space B is connected and locally path-connected, the map \(\varphi\) is surjective, continuous, and open (III, p. 262, prop. 10). The strict homotopy relation is therefore open (TG, I, p. 32, prop. 3).

In the general case, let E be a simply path-connected covering of B, nonempty if \(B \neq \varnothing\) . This covering is surjective because B is connected (I, p. 74, prop. 4). Let p denote the projection of the B-space E and let \(\widetilde{p} \colon \Lambda(\mathrm{E}) \to \Lambda(\mathrm{B})\) be the map which, to a path c in E, associates the path \(p \circ c\) . Equipped with the map \(\widetilde{p}\) , \(\Lambda(\mathrm{E})\) is a covering of \(\Lambda(\mathrm{B})\) (III, p. 302, cor. 1 of prop. 3). For two paths in B to be strictly homotopic, it is necessary and sufficient that they be the images under the map \(\widetilde{p}\) of two strictly homotopic paths in E (III, p. 302, prop. 4). Let U be an open subset of the space \(\Lambda(B)\) . The set \((\widetilde{p})(U)\) is open in \(\Lambda(E)\) ; by the first part of the proof, the set \(U'\) saturated with respect to \((\widetilde{p})(U)\) for the strict homotopy relation is open in \(\Lambda(E)\) . Since \(\Lambda(E)\) is a covering of \(\Lambda(B)\) the set \(\widetilde{p}(U')\) is open in \(\Lambda(B)\) . This proves the proposition, since \(\widetilde{p}(U')\) is the saturation of U for the strict homotopy relation.

\subsection{Examples}

\subitemhead{Locally contractible spaces}

DEFINITION 3. --- A topological space B is said to be locally contractible if every point b of B has a neighborhood V such that the pointed space (V, b) is contractible (III, p. 234, example 3).

PROPOSITION 5. --- A locally contractible topological space is loop-disentanglable.

Let B be a locally contractible space. Let b be a point of B and let V be a neighborhood of b such that \((V, b)\) is a pointed contractible space. The space V is homotopy equivalent to a point, hence \(\pi_{1}(V, b) = \{\varepsilon_{b}\}\) (III, p. 296, corollary 3).

To prove the proposition, it remains to show that the space B is locally path-connected. Let \(\sigma: V \times I \to V\) be a pointed homotopy at \(b\) relating the constant map with image \(b\) to the map \(Id_V\) . For every neighborhood W of \(b\) contained in V, the set \(\sigma(W \times I)\) is a neighbourhood of \(b\) since it contains \(W = \sigma(W \times \{1\})\) , and it is path-connected because for every \((a, s) \in W \times I\) , the map \(t \mapsto \sigma(a, ts)\) is a path connecting \(b = \sigma(a, 0)\) to \(\sigma(a, s)\) in \(\sigma(W \times I)\) . Let us prove that the sets of the form \(\sigma(W \times I)\) form a fundamental system of neighborhoods of \(b\) . Let \(V'\) be an open neighborhood of \(b\) in B; then, \(\sigma^{-1}(V')\) is an open neighborhood of \(\{b\} \times I\) in \(V \times I\) . Since the space I is compact, the projection \(\mathrm{pr}_1: V \times I \to V\) is proper (TG, I, p. 77, cor. 5) and \(\mathrm{pr}_1(\complement\,\sigma^{-1}(V'))\) is a closed subset of V not containing \(b\) . Its complement W is thus an open neighborhood of \(b\) in V such that \(\mathrm{W} \times \mathbf{I} \subset \sigma^{-1}(\mathrm{V}')\) , whence \(\sigma(\mathrm{W} \times \mathbf{I}) \subset \mathrm{V}'\) .

Every open subset of a numerical space or of a projective space, real or complex, of dimension n (cf. chapters VI and VIII) is loop-disentanglable, as are the Euclidean spheres \(S_{n}\) (TG, VI, p. 11, prop. 4). More generally, every topological manifold (VAR R, second part, p. 7, Notations and Conventions) is loop-disentanglable. Indeed, these spaces are locally contractible.

\subitemhead{Poincaré group of the circle}

The topological space R is homotopic to a point (III, p. 234, example 4), hence simply path-connected (IV, p. 340). The canonical surjection p of R onto T = R/Z makes it a principal covering with group Z (I, p. 100, example 4). The topological space T is loop-disentanglable and (R, 0) is a universal covering of (T, 0). From this we deduce a canonical isomorphism of groups \(\pi_{1}(\mathbf{T}, 0) \to \mathbf{Z}\) : it maps the class of the loop \(t \mapsto p(nt)\) at 0 to the element n of Z. Taking into account example 6 (I, p. 101), we deduce the following proposition.

Proposition 6. — The map \(p \colon x \mapsto e^{2\pi ix}\) from \(\mathbf{R}\) onto \(\mathbf{S}_1\) makes \((\mathbf{R}, 0)\) a universal covering of \((\mathbf{S}_1, 1)\) . The group \(\pi_1(\mathbf{S}_1, 1)\) is isomorphic to \(\mathbf{Z}\) ; the class of the loop \(t \mapsto e^{2\pi it}\) is a generator. For every integer \(n > 0\) , the map \(z \mapsto z^n\) makes \(\mathbf{S}_1\) a Galois covering with group \(\mathbf{Z}/n\mathbf{Z}\) of \(\mathbf{S}_1\) . Every connected and nonempty covering of \(\mathbf{S}_1\) is isomorphic to one of these coverings.

Only the last assertion remains to be proved. The fiber \(E_{1}\) over 1 of a connected and nonempty covering E of \(\mathbf{S}_{1}\) is equipped with a transitive action of the group \(\pi_{1}(\mathbf{S}_{1},1)\) . Since the subgroups of Z are the sets nZ, for \(n \geqslant 0\) , one observes that \(E_{1}\) is isomorphic to one of the homogeneous sets associated with the coverings described previously. Since \(\mathbf{S}_{1}\) is connected and locally path-connected, E is isomorphic to one of these coverings.

COROLLARY. --- The map \(z \mapsto e^{z}\) from C to \(C^{*}\) makes \((\mathbf{C},0)\) a universal covering of \((\mathbf{C}^{*},1)\) . The group \(\pi_{1}(\mathbf{C}^{*},1)\) is isomorphic to Z; the class of the loop \(t \mapsto e^{2\pi it}\) is a generator. For every integer n \textgreater{} 0, the map \(z \mapsto z^{n}\) makes \(C^{*}\) a Galois covering with group Z/nZ of \(C^{*}\) . Every connected and nonempty covering of \(C^{*}\) is isomorphic to one of these coverings.

The map \(z \mapsto (|z|, z/|z|)\) is a homeomorphism of the space \(\mathbf{C}^*\) onto the space \(\mathbf{R}_+^* \times \mathbf{S}^1\) (TG, VI, p. 10, prop. 3); in particular, \(\mathbf{C}^*\) is path-connected. Since the map \(x \mapsto e^x\) from \(\mathbf{R}\) onto \(\mathbf{R}_+^*\) is a homeomorphism, it follows from Proposition 6 that the map \((x, y) \mapsto e^x e^{2\pi iy}\) makes \(\mathbf{R}^2\) a covering of \(\mathbf{C}^*\) . By identifying \(\mathbf{R}^2\) with \(\mathbf{C}\) by the mapping \((x, y) \mapsto x + iy\) , we conclude that the space \(\mathbf{C}\) , equipped with the map \(z \mapsto e^z\) , is a covering of \(\mathbf{C}^*\) . Since the space \(\mathbf{C}\) is connected and simply path-connected, the pointed covering \((\mathbf{C}, 0)\) is a universal covering of the pointed space \((\mathbf{C}^*, 1)\) .

It also follows from the corollary, III, p. 297, that the Poincaré group of \(\mathbf{C}^*\) at 1 is isomorphic to \(\mathbf{Z}\) and that the class \(\gamma\) of the loop \(t \mapsto e^{2\pi it}\) is a generator.

  Let n be an integer \textgreater{} 0. The map \(x \mapsto x^{n}\) from \(R_{+}^{*}\) onto itself is a homeomorphism; we deduce as before that the map \(z \mapsto z^{n}\) makes \(C^{*}\) a covering of degree n of \(C^{*}\) . It is Galois with group Z/nZ. The last assertion is proved as in the proof of Prop. 6.

For every integer \(n \geqslant 0\) , the torus \((\mathbf{S}_{1})^{n}\) is a loop-disentanglable space (IV, p. 341, prop. 1) and its Poincaré group at every point is isomorphic to \(\mathbf{Z}^{n}\) (III, p. 297, prop. 4).

We will generalize these results in Section 3, which is devoted to coverings of topological groups.

\subitemhead{Real projective spaces}

Let n be an integer \textgreater{} 0. The spaces \(S_{n}\) and \(\mathbf{P}_{n}(\mathbf{R})\) are connected, nonempty, and the canonical map (TG, VI, p. 13) \(\varphi\colon\mathbf{S}_{n}\to\mathbf{P}_{n}(\mathbf{R})\) makes \(S_{n}\) a principal covering of \(\mathbf{P}_{n}(\mathbf{R})\) of the group \(\{+1,-1\}\) (I, p. 99, example 1). For \(n\geqslant2\) the sphere \(S_{n}\) is simply connected (I, p. 127, example 3) and loop-disentanglable, hence simply path-connected (IV, p. 344, corollary 1). For n=1, the map \(z\mapsto z^{2}\) from \(S_{1}\) to \(S_{1}\) defines an equivalence relation in \(S_{1}\) whose classes are the pairs of opposite points of \(S_{1}\) . The map \(\varphi\colon\mathbf{S}_{1}\to\mathbf{P}_{1}(\mathbf{R})\) defined by passing to the quotient defines a continuous bijection \(\psi\colon\mathbf{S}_{1}\to\mathbf{P}_{1}(\mathbf{R})\) such that \(\psi(z^{2})=\varphi(z)\) . Since \(S_{1}\) is a compact space, the map \(\psi\) is a homeomorphism (TG, I, p. 63, cor. 2).

PROPOSITION 7. --- The map \(x \mapsto \varphi(e^{2\pi ix})\) from \(\mathbf{R}\) onto \(\mathbf{P}_{1}(\mathbf{R})\) makes \((\mathbf{R},0)\) a universal covering of \((\mathbf{P}_{1}(\mathbf{R}),\varphi(1))\) . Let c: \(\mathbf{I} \to \mathbf{S}_{1}\) be the path \(t \mapsto e^{\pi it}\) ; the class of \(\varphi(c)\) is a generator of the group \(\pi_1(\mathbf{P}_{1}(\mathbf{R}), \varphi(1))\) which is isomorphic to \(\mathbf{Z}\) .

For every integer \(n \geqslant 2\) and every point x of \(S_{n}\) , the canonical map \(\varphi: \mathbf{S}_{n} \to \mathbf{P}_{n}(\mathbf{R})\) makes \((\mathbf{S}_{n}, x)\) a universal covering of \((\mathbf{P}_{n}(\mathbf{R}), \varphi(x))\) . For every path c in \(S_{n}\) joining x to -x, the class of \(\varphi \circ c\) generates the group \(\pi_{1}(\mathbf{P}_{n}(\mathbf{R}), \varphi(x))\) which is isomorphic to Z/2Z.

\subitemhead{Quotient of a space by a proper action of a discrete group}

Lemma 1. --- Let X be a topological space, let G be a discrete group acting properly on X, and let p be the canonical projection from X onto X/G.

Let \(x\) be a point of \(X\) , let \(K_x\) be its fixator in \(G\) . Let \(U_1\) be a neighborhood of \(x\) in \(X\) ; there exists a neighborhood \(U\) of \(x\) in \(X\) stable under \(K_x\) and contained in \(U_1\) such that \(g \cdot U \cap U = \varnothing\) if \(g \notin K_x\) and such that the canonical map \(p\) induces a homeomorphism of \(U / K_x\) onto a neighborhood \(V\) of \(p(x)\) in \(X / G\) .

By prop. 8 of TG, III, p. 32, there exists a neighborhood \(U_{2}\) of x in X, stable under \(K_{x}\) , such that \(g \cdot U_{2} \cap U_{2} = \varnothing\) if \(g \notin K_{x}\) and such that the canonical map p induces a homeomorphism from \(U_{2}/K_{x}\) onto a neighborhood of \(p(x)\) in X/G.

As the canonical map \(U_{2} \rightarrow U_{2}/K_{x}\) is closed (TG, III, p. 28, prop. 2), there exists an open neighborhood U of x contained in \(U_{1} \cap U_{2}\) which is stable under \(K_{x}\) . The equivalence relation in \(U_{2}\) defined by \(K_{x}\) is also open (TG, III, p. 10, lemma 2). It then follows from TG, I, p. 32, prop. 4, that the canonical mapping induces a homeomorphism from \(U/K_{x}\) onto an open neighborhood of \(p(x)\) in X/G, whence the lemma.

PROPOSITION 8. --- Let X be a topological space and let G be a discrete group acting properly on X. Denote by p the canonical projection from X to X/G.

\begin{enumerate}
\def\labelenumi{\alph{enumi})}
\item
  Suppose that X is path-connected and that the group G is generated by the stabilizers of the points of X. Then the canonical homomorphism \(\pi_{1}(X,x)\to\pi_{1}(X/G,p(x))\) is surjective for every point x in X. In particular, X/G is simply path-connected if X is.
\item
  If the space X is loop-disentanglable, then the space X/G is loop-disentanglable.
\item
  The set N of elements \(g \in G\) for which there exists a path \(c_{g} \colon I \to X\) such that \(c_{g}(0) = x\) , \(c_{g}(1) = g \cdot x\) and such that the loop \(p \circ c_{g}\) in X/G is strictly homotopic to a constant loop is a subgroup of G. If \(g \in G\) and if there exists a point y of X such that \(g \cdot y = y\) , choose a path c: I → X joining x to y; then the path \(c' = c * (g \cdot c)\) satisfies \(c'(0) = x\) , \(c'(1) = g \cdot x\) and \([p \circ c'] = [p \circ c][p \circ (g \cdot c)]^{-1} = e_{p(x)}\) , therefore \(p \circ c'\) is strictly homotopic to a constant loop. Since G is generated by the fixators of the points of X, it follows that N = G.
\end{enumerate}

Let \(c\) be a loop in X/G at \(p(x)\) . By Theorem 4 of III, p. 287, there exists a path \(\widetilde{c} \colon \mathbf{I} \to \mathrm{X}\) lifting \(c\) such that \(\widetilde{c}(0) = x\) . Since \(p(\widetilde{c}(1)) = c(1) = p(x)\) , there exists \(g \in \mathrm{G}\) such that \(\widetilde{c}(1) = g \cdot x\) . Choose a path \(c_g \colon \mathbf{I} \to \mathrm{X}\) connecting \(x\) to \(g \cdot x\) and such that \(p \circ c_g\) is strictly homotopic to a constant loop. Then the path \(c' = \widetilde{c} * \overline{c_g}\) is a loop at \(x\) in X such that \([p \circ c'] = [c]\) . This shows that the homomorphism \(\pi_1(\mathrm{X}, x) \to \pi_1(\mathrm{X/G}, p(x))\) is surjective. The other assertion follows immediately.

Let us now prove assertion \(b\) ). The space X/G is locally path-connected (III, p. 261, prop. 8).

Let \(x\) be a point of X and let \(K_x\) be its stabilizer in G. Let \(U_1\) be an open neighborhood of \(x\) such that the image of the canonical homomorphism \(\pi_1(U_1, x) \to \pi_1(X, x)\) is reduced to the identity element. By lemma 1 (IV, p. 349), there exists a neighborhood U of \(x\) in X, contained in \(U_1\), stable under \(K_x\) , such that \(g \cdot U \cap U = \varnothing\) if \(g \notin K_x\) and such that the canonical map \(p\) induces a homeomorphism of \(U/K_x\) onto a neighborhood V of \(p(x)\) in X/G.

Let \(c\) be a loop at \(p(x)\) in V; by theorem 4 (III, p. 287), applied to the topological space U and the group \(K_x\) , there exists a path \(\widetilde{c} \colon \mathbf{I} \to \mathbf{U}\) such that \(\widetilde{c}(0) = x\) and which lifts \(c\) . One necessarily has \(\widetilde{c}(1) = x\) , so that \(\widetilde{c}\) is a loop at \(x\) in U. By hypothesis, \(\widetilde{c}\) is strictly homotopic to a constant loop in X; consequently, the same holds for \(c\) and the canonical homomorphism \(\pi_1(V, p(x)) \to \pi_1(X/G, p(x))\) is trivial. This shows that X/G is loop-disentanglable if X is.

\begin{enumerate}
\def\labelenumi{\arabic{enumi})}
\setcounter{enumi}{4}
\tightlist
\item
  In the Euclidean plane \(R^{2}\) , let P be the topological space which is the union of the circles with center \((1/n,0)\) passing through the origin (for \(n \in N^{*}\) ). The space P is compact, connected, and locally path-connected, but is not loop-disentanglable. The group \(\pi_{1}(P,0)\) , equipped with the admissible topology, is Hausdorff and non-discrete (III, p. 337, exerc. 7).
\end{enumerate}

Remark 1. --- Theorem 4 of III, p. 287 and prop. 8 (IV, p. 349) remain valid under the more general hypothesis that G is a finite-dimensional Lie group over R acting properly on X. For more details, cf. D. MONTGOMERY and C. T. YANG, “The existence of a slice”, Annals of mathematics 65 (1957), p. 108--116; R. PALAIS, “On the existence of slices for actions of non-compact Lie groups”, Annals of mathematics 73 (1961), p. 295--323; and G. BREDON, Introduction to compact transformation groups, Academic Press, 1972.

\section{POINCARÉ GROUPS OF LOOP-DISENTANGLABLE SPACES}

\subsection{\texorpdfstring{Properties of homomorphisms \(\pi_{1}(f,a)\)}{Properties of homomorphisms pi\_1(f,a)}}

Let A and B be topological spaces, let \((\mathrm{E},p)\) be a covering of A, let \((\mathrm{E}^{\prime},p^{\prime})\) be a covering of B, let \(f\colon A\to B\) and \(g\colon E\to E^{\prime}\) be continuous maps such that \(p^{\prime}\circ g=f\circ p\) . Let a be a point of A and let \(b=f(a)\) . For all \(\gamma\in\pi_{1}(\mathrm{A},a)\) and every \(x\in E_{a}\) , we have \(g(x\cdot\gamma)=g(x)\cdot f_{*}(\gamma)\) , according to relation (3), III, p. 304. If the diagram

\[
\begin{array}{c} \mathrm{E} \xrightarrow {g} \mathrm{E} ^ {\prime} \\ \Big \downarrow p \\ \mathrm{A} \xrightarrow {f} \mathrm{B} \end{array}
\]

is a Cartesian square, the map g induces a bijection from \(E_{a}\) onto \(E'_{b}\) . The action of \(\pi_{1}(A,a)\) on the fiber \(E_{a}\) is then the composite of the action of the group \(\pi_{1}(B,b)\) on \(E'_{b}\) and the homomorphism \(\pi_{1}(f,a)\) from \(\pi_{1}(A,a)\) into \(\pi_{1}(B,b)\) .

PROPOSITION 1. --- Let A be a loop-disentanglable topological space, B a locally path-connected topological space, and let \(f\colon\mathrm{A}\to\mathrm{B}\) be a continuous map. Let a be a point of A and \(b=f(a)\) . Suppose that every covering of A is isomorphic to the inverse image under \(f\) of a covering of B. Then the homomorphism \(\pi_{1}(f,a)\colon\pi_{1}(\mathrm{A},a)\to\pi_{1}(\mathrm{B},b)\) is injective.

Since the space A is loop-disentanglable, there exists a covering E of A whose fiber \(E_{a}\) is a principal homogeneous \(\pi_{1}(A,a)\) -set (IV, p. 342, theorem 1). By hypothesis, this action is induced via the homomorphism \(\pi_{1}(f,a)\) by an action of \(\pi_{1}(B,b)\) on \(E_{a}\) . This implies the injectivity of the homomorphism \(\pi_{1}(f,a)\) .

PROPOSITION 2. --- Let A be a connected and locally path-connected topological space, B a loop-disentanglable topological space, and let \(f\colon\mathrm{A}\to\mathrm{B}\) be a continuous map. Let a be a point of A and \(b=f(a)\) . The following assertions are equivalent:

\begin{enumerate}
\def\labelenumi{(\roman{enumi})}
\item
  The homomorphism \(\pi_1(f,a)\colon \pi_1(A,a)\to \pi_1(B,b)\) is surjective.
\item
  For every pair (E, E') of coverings of B, the map
\end{enumerate}

\[
f ^ {*} \colon \mathscr {C} _ {\mathrm{B}} (\mathrm{E}; \mathrm{E} ^ {\prime}) \rightarrow \mathscr {C} _ {\mathrm{A}} (\mathrm{A} \times_ {\mathrm{B}} \mathrm{E}; \mathrm{A} \times_ {\mathrm{B}} \mathrm{E} ^ {\prime})
\]

which, to a B-morphism \(g\colon \mathrm{E}\to \mathrm{E}^{\prime}\) , associates the A-morphism

\[
f ^ {*} (g) \colon \mathrm{A} \times_ {\mathrm{B}} \mathrm{E} \rightarrow \mathrm{A} \times_ {\mathrm{B}} \mathrm{E} ^ {\prime}, \quad (x, y) \mapsto (x, g (y))
\]

is bijective.

By means of the homomorphism \(\pi_1(f, a)\) , every \(\pi_1(B, b)\)-set is equipped with a structure of \(\pi_1(A, a)\) -set. Condition (i) is then equivalent to the following condition (i'):

(i') For every pair \((\mathrm{F},\mathrm{F}^{\prime})\) of \(\pi_{1}(\mathrm{B},b)\) -sets, every \(\pi_{1}(\mathrm{A},a)\) -morphism from F to \(F^{\prime}\) is a \(\pi_{1}(\mathrm{B},b)\) -morphism. Indeed, if the homomorphism \(\pi_{1}(f,a)\) is surjective, every \(\pi_{1}(\mathrm{A},a)\) -morphism of \(\pi_{1}(\mathrm{B},b)\) -sets is a \(\pi_{1}(\mathrm{B},b)\) -morphism. Conversely, take a \(\pi_{1}(\mathrm{B},b)\) -set F reduced to a point and set \(\mathrm{F}^{\prime}=\pi_{1}(\mathrm{B},b)/f_{*}(\pi_{1}(\mathrm{A},a))\) . The map from F to \(F^{\prime}\) whose image is \(f_{*}(\pi_{1}(\mathrm{A},a))\) is a \(\pi_{1}(\mathrm{A},a)\) -morphism but is not a \(\pi_{1}(\mathrm{B},b)\) -morphism if \(F^{\prime}\) is not reduced to a point, that is, if \(\pi_{1}(f,a)\) is not surjective.

Since the space B is loop-disentanglable, every \(\pi_{1}(B,b)\) -set is isomorphic to the \(\pi_{1}(B,b)\) -set \(E_{b}\) , where E is a covering of B (IV, p. 344, remark 1). The equivalence of (i') and (ii) therefore follows from prop. 2 of III, p. 310.

PROPOSITION 3. --- Let B be a loop-disentanglable topological space and let A be a connected and locally path-connected subspace of B (for example an open and connected subset). Let a be a point of A. The following properties are equivalent:

\begin{enumerate}
\def\labelenumi{(\roman{enumi})}
\item
  Every covering of B is trivializable over A.
\item
  The image of the homomorphism from \(\pi_{1}(\mathrm{A},a)\) into \(\pi_{1}(\mathrm{B},a)\) deduced from the canonical injection is reduced to the identity element.
\end{enumerate}

The implication (ii)⇒(i) follows from corollary 3 (III, p. 309).

Conversely, suppose that every covering of B is trivializable over A. This is in particular the case for the simply path-connected covering \((\lambda_{a}(\mathrm{B}),\varepsilon_{\mathrm{B}})\) (IV, p. 342, th. 1), so that the homomorphism \(\pi_{1}(\mathrm{A},a)\to\pi_{1}(\mathrm{B},a)\) is trivial (III, p. 309, cor. 3 of prop. 1).

\subsection{Relatively connected maps}

DEFINITION 1. --- Let X and Y be topological spaces and let f be a continuous map from X to Y. The map f is said to be relatively connected if every point of Y has a fundamental system of neighborhoods consisting of sets V such that \(f^{-1}(V)\) is connected and nonempty.

Let \(f \colon X \to Y\) be a continuous map; for \(f\) to be relatively connected, it is necessary and sufficient that every point of \(Y\) has a neighborhood \(V\) such that the map \(f_V \colon f^{-1}(V) \to V\) deduced from \(f\) be relatively connected.

Let X and Y be topological spaces, and let \(f: X \to Y\) be a relatively connected continuous map. The image of f is dense in Y. For every open subset V of Y, the map \(f_{\mathrm{V}}: f^{-1}(\mathrm{V}) \to \mathrm{V}\) is relatively connected.

PROPOSITION 4. --- Let X and Y be topological spaces and let \(f: X \to Y\) be a relatively connected continuous map.

\begin{enumerate}
\def\labelenumi{\alph{enumi})}
\item
  For every connected component U of X, \(f(\overline{\mathrm{U}})\) is a connected, open, and closed subset of Y, and we have \(f^{-1}(\overline{f(\mathrm{U})}) = \mathrm{U}\) .
\item
  For every connected component V of Y, there exists a connected component U of X such that \(V = \overline{f(U)}\) . The map from U to V induced by f by passing to subsets is relatively connected.
\item
  The connected components of X (resp. of Y) are open and closed.
\item
Let V be the set of \(y \in Y\) which possess a neighborhood W such that \(f^{-1}(W)\) is connected and meets U; it is an open subset of Y. It contains \(\overline{f(U)}\) since f is relatively connected. Conversely, let \(y \in V\); let W be a neighborhood of \(y\) such that \(f^{-1}(W)\) is connected and meets U. For every neighborhood \(W'\) of \(y\) such that \(f^{-1}(W')\) is connected, \(f^{-1}(W \cap W')\) is not empty, therefore \(f^{-1}(W \cup W')\) is connected (TG, I, p. 81, prop. 2). By hypothesis, \(f^{-1}(W \cup W')\) meets U; one therefore has \(f^{-1}(W \cup W') \subset U\) and, a fortiori, \(W' \cap f(U) \neq \varnothing\). This proves that \(y \in \overline{f(U)}\); therefore, \(V = \overline{f(U)}\). In particular, the set \(\overline{f(U)}\) is open and closed in Y; moreover, prop. 1 (TG, I, p. 81) implies that it is connected.
\end{enumerate}

The preceding arguments also show that \(f^{-1}(\overline{f(\mathrm{U})}) \subset \mathrm{U}\) , whence the equality \(f^{-1}(\overline{f(\mathrm{U})}) = \mathrm{U}\) , the other inclusion being evident. In particular, U is open and closed in X.

\begin{enumerate}
\def\labelenumi{\alph{enumi})}
\setcounter{enumi}{1}
\tightlist
\item
Let V be the connected component of a point y of Y, let W be a neighborhood of y such that \(f^{-1}(W)\) is connected and nonempty, and let U be the connected component of X that contains \(f^{-1}(W)\). Since \(y \in \overline{f(U)}\), it follows from a) and from the definition of the connected component of y that \(\mathrm{V} = \overline{f(\mathrm{U})}\). Consequently, V is open and closed in Y.
\end{enumerate}

COROLLARY 1. --- Let X and Y be topological spaces and let \(f: X \to Y\) be a continuous and relatively connected map. By passing to connected components, the map \(f\) induces a bijection from the set of connected components of X onto the set of connected components of Y.

COROLLARY 2. --- Let X and Y be topological spaces and let \(f: X \to Y\) be a continuous map. For \(f\) to be relatively connected, it is necessary and sufficient that the following three properties be satisfied:

\begin{enumerate}
\def\labelenumi{\alph{enumi})}
\item
  The image f(X) is dense in Y;
\item
  The space Y is locally connected;
\item
  For every open and connected set V in Y, the set \(f(\mathrm{V})\) is connected.
\end{enumerate}

Suppose that the map f is relatively connected. The density of \(f(\mathrm{X})\) in Y follows from the definition of a relatively connected map. Let V be an open subset of Y; the map \(f_{\mathrm{V}}: f^{-1}(\mathrm{V}) \to \mathrm{V}\) is relatively connected. By prop. 4, the connected components of V are open and closed in V. It follows that Y is locally connected (TG, I, p. 85, prop. 11). To prove assertion c), it suffices to prove that X is connected if Y is. By the lemma, there exists a connected component U of X such that \(Y = \overline{f(U)}\) and \(f^{-1}(\overline{f(U)}) = U\) , whence U = X and X is connected.

Conversely, suppose that the conditions \(a\), \(b\), \(c\) are satisfied and let us show that f is relatively connected. Let \(y\) be a point of Y; since Y is locally connected, \(y\) admits a fundamental system of connected open neighborhoods. If W is such a neighborhood, the conditions \(c\) and \(a\) imply that \(f^{-1}(W)\) is connected and nonempty. Consequently, the map f is relatively connected.

COROLLARY 3. --- Let X and Y be topological spaces and let \(f\colon X\to Y\) be a continuous and relatively connected map. Let F be a set and let \(g\colon X\to F\) be a locally constant map. There exists a unique locally constant map \(h\colon Y\to F\) such that \(g=h\circ f\) .

The restriction of g to every connected component of X is constant. By Corollary 1, there exists a map \(h: Y \to F\) , constant on each connected component of Y, such that \(g = h \circ f\) . The map h is locally constant, because the connected components of Y are open. The uniqueness of such a map follows from the fact that \(f(X)\) is dense in Y.

Examples. --- 1) Let X be a topological space and let R be an equivalence relation on X. Denote by Y the quotient topological space X/R and by \(f\colon X \to Y\) the canonical map. Suppose that the equivalence classes of R are connected. Then, for every open and connected subset V of Y, the set \(f^{-1}(V)\) is connected (TG, I, p. 23, corollary 1 and p. 82, proposition 7). If the space Y is locally connected, the map f is thus relatively connected.

\begin{enumerate}
\def\labelenumi{\arabic{enumi})}
\setcounter{enumi}{1}
\tightlist
\item
  Let Y be a locally path-connected topological space and let X be an open subset of Y. The space \(\mathcal{C}_{c}(\mathbf{I};\mathbf{X})\) of paths in X is identified with a subspace of the space \(\mathcal{C}_{c}(\mathbf{I};\mathbf{Y})\) of paths in Y; suppose it is dense. The canonical injection of X into Y is then a relatively connected map.
\end{enumerate}

It suffices indeed to show that, for every nonempty connected open subset V of Y, the set \(V \cap X\) is connected and nonempty. The space \(\mathcal{C}_{c}(\mathbf{I};\mathrm{V})\) is a nonempty open set of \(\mathcal{C}_{c}(\mathbf{I};\mathrm{Y})\) ; it meets \(\mathcal{C}_{c}(\mathbf{I};\mathrm{X})\) , which proves that \(V \cap X \neq \varnothing\) . Let x and \(x'\) be points of \(V \cap X\) . Since \(V \cap X\) is locally path-connected, the points x and \(x'\) have open neighborhoods U and \(U'\) path-connected and contained in \(V \cap X\) . Since V is path-connected (III, p. 260, prop. 7), there exists a path in V joining x to \(x'\) . By the density hypothesis and the definition of the compact-convergence topology, there exists a path in \(V \cap X\) joining a point of U to a point of \(U'\) . There then exists a path joining x to \(x'\) in \(V \cap X\) , which proves that the set \(V \cap X\) is connected.

PROPOSITION 5. --- Let X and Y be topological spaces and let \(f: X \to Y\) be a continuous and relatively connected map. For every pair \((T, T')\) of coverings of Y, the map \(f^{*}: \mathcal{C}_{Y}(T; T') \to \mathcal{C}_{X}(X \times_{Y} T; X \times_{Y} T')\) is bijective.

Let \(\mathcal{F}\) be the sheaf on X of X-morphisms from \(X \times_{Y} T\) into \(X \times_{Y} T'\) and let \(\mathcal{G}\) be the sheaf on Y of Y-morphisms from T into \(T'\) (I, p. 45, example 4). For every open set U of Y, set \(\varphi_{\mathrm{U}} = (f_{\mathrm{U}})^{*}: \mathcal{G}(\mathrm{U}) \to \mathcal{F}(f^{-1}(\mathrm{U}))\) . The maps \(\varphi_{\mathrm{U}}\) define a morphism of sheaves \(\varphi: \mathcal{G} \to \varphi_{*}(\mathcal{F})\) and it suffices to prove that \(\varphi\) is an isomorphism of sheaves.

Since Y is locally connected (IV, p. 354, cor. 2), the connected open sets over which T and T' are trivializable form a base of the topology of Y. By corollary 2 of I, p. 55, it suffices to prove that for such an open set U, the map \(\varphi_{\mathrm{U}}\) is bijective, which allows us to assume that Y is connected and that the coverings T and T' are the trivial coverings \(Y \times F\) and \(Y \times F'\) where F and F' are sets equipped with the discrete topology. The map \((x, (y, t)) \mapsto (x, t)\) identifies the X-space \(X \times_{Y} (Y \times F)\) with \(X \times F\) (resp. the X-space \(X \times_{Y} (Y \times F')\) with \(X \times F'\)). Since the space X is connected (IV, p. 354, cor. 2), the sets \(\mathcal{C}_{\mathrm{Y}}(Y \times F; Y \times F')\) and \(\mathcal{C}_{\mathrm{X}}(X \times F; X \times F')\) are both identified with the set \(\mathcal{F}(F;F')\) of maps from F into F', and the map \(f^{*}\) is identified with the identity map of \(\mathcal{F}(F;F')\). This concludes the proof.

COROLLARY 1. --- Let X and Y be topological spaces and let \(f: X \to Y\) be a relatively connected continuous map. Let T and T' be coverings of Y. If the coverings \(X \times_Y T\) and \(X \times_Y T'\) of X are isomorphic, then the coverings T and T' are isomorphic.

Indeed, let \(h \colon X \times_{Y} T \to X \times_{Y} T'\) and \(h' \colon X \times_{Y} T' \to X \times_{Y} T\) be X-isomorphisms reciprocal to each other. By Proposition 5, there exist Y-morphisms \(g \colon \mathrm{T} \to \mathrm{T}'\) and \(g' \colon \mathrm{T}' \to \mathrm{T}\) such that \(f^{*}(g) = h\) and \(f^{*}(g') = h'\) . One then has \(f^{*}(g' \circ g) = f^{*}(g') \circ f^{*}(g) = \mathrm{Id}_{\mathrm{X} \times_{\mathrm{Y}} \mathrm{T}}\) , therefore \(g' \circ g = \mathrm{Id}_{\mathrm{T}}\) , because the map \(f^{*}\) is injective. Similarly, \(g \circ g' = \mathrm{Id}_{\mathrm{T}'}\) . The coverings T and \(\mathrm{T}'\) are therefore isomorphic.

COROLLARY 2. --- Let X and Y be topological spaces and let \(f: X \to Y\) be a continuous and relatively connected map. If the space X is simply connected, then so is the space Y.

By Corollary 1, every covering of Y is indeed trivializable.

COROLLARY 3. --- Let X be a locally path-connected topological space, Y a loop-disentanglable topological space, and let \(f: X \to Y\) be a relatively connected continuous map. For every point \(x\) of X, the homomorphism \(\pi_1(f, x): \pi_1(X, x) \to \pi_1(Y, f(x))\) is surjective.

By IV, p. 353, prop. 4, we may assume that the spaces X and Y are connected. The corollary then follows from proposition 5 and proposition 2 of IV, p. 352.

PROPOSITION 6. --- Let X and Y be topological spaces, and let \(f: \mathrm{X} \to \mathrm{Y}\) be a continuous and relatively connected map. Suppose that every point of Y has an open neighborhood V whose inverse image \(f^{-1}(V)\) is simply connected. Then, for every covering Z of X, there exists a covering T of Y such that \(X \times_{Y} T\) is X-isomorphic to Z.

Let \(\mathcal{U}\) be the set of open subsets of Y whose inverse image in X is simply connected. For every \(V \in \mathcal{U}\), the covering \(f^{-1}(V) \times_{X} Z\) of \(f^{-1}(V)\) is trivializable; there thus exists a discrete space \(F_V\) and an isomorphism of coverings \(g_V: f^{-1}(V) \times_{X} Z \to f^{-1}(V) \times F_V\) . For every pair \((V, V')\) of open sets belonging to \(\mathcal{U}\), the map \(f_{V \cap V'}: f^{-1}(V \cap V') \to V \cap V'\) is relatively connected. By Proposition 5 of IV, p. 356, there exists a unique isomorphism of coverings of \(V \cap V'\), \(h_{V',V}: (V \cap V') \times F_V \to (V \cap V') \times F_{V'}\) such that one has \(f^*(h_{V',V})(x,t) = g_{V'}(g_V^{-1}(x,t))\) for all \(x \in V \cap V'\) and every \(t \in F_V\). If V, \(V'\), \(V''\) are elements of \(\mathcal{U}\), we have \(h_{V'',V}(x,t) = h_{V'',V'}(h_{V',V}(x,t))\) for all \(x \in V \cap V' \cap V''\) and every \(t \in F_V\). There then exists a unique Y-space T and, for every \(V \in \mathcal{U}\), an isomorphism \(h_V: T_V \to V \times F_V\) such that one has \(h_{V',V}(x,t) = h_{V'} \circ h_V^{-1}(x,t)\) for every pair \((V, V')\) of open sets belonging to \(\mathcal{U}\), all \(x \in V \cap V'\) and every \(t \in F_V\) (cf. TG, I, p. 16).

The space T is in particular a covering of Y. Moreover, there exists a unique map from \(X \times_{Y} T\) over Z whose restriction to \(f^{-1}(V) \times_{Y} T\) is given by \(g_{V}^{-1} \circ f^{*}(h_{V})\) and it is an isomorphism of X-spaces, whence the proposition.

COROLLARY. --- Let X and Y be loop-disentanglable topological spaces and let \(f: X \to Y\) be a continuous and relatively connected map. Suppose that every point of Y has a neighborhood V whose inverse image \(f^{-1}(V)\) is simply connected. Then, for every point x in X, the homomorphism \(\pi_{1}(f,x): \pi_{1}(X,x) \to \pi_{1}(Y,f(x))\) is bijective.

One may assume the spaces X and Y are connected (IV, p. 353, prop. 4). By Corollary 3 (IV, p. 357), the homomorphism \(\pi_1(f,x)\) is surjective. Proposition 6 and Proposition 1 of IV, p. 351 imply that it is injective.

Remark. --- Let Y be a topological space, and let X, X', and Y' be subspaces of Y such that X ⊂ X' ⊂ Y' ⊂ Y. Suppose that the canonical injection of X into Y is relatively connected. For every connected open subset V of Y, the set V ∩ X is connected and dense in V (IV, p. 354, cor. 2); consequently, the set V ∩ X' is connected (TG, I, p. 81, prop. 1). This proves that the canonical injection of X' into Y' is relatively connected.

Let V be an open subset of Y; by the preceding, the canonical injection of \(V \cap X\) into \(V \cap X'\) is relatively connected. If the set \(V \cap X\) is simply connected, \(V \cap X'\) is also, by virtue of Corollary 2 of IV, p. 357. Consequently, if the canonical injection of X into Y satisfies the hypotheses of Proposition 6, then so does the canonical injection of \(X'\) into \(Y'\) .

Example. --- Let Y be a locally finite-dimensional differentiable manifold and let Z be a closed submanifold of Y (VAR, R, 5.8.3). Set X = Y - Z.

\begin{enumerate}
\def\labelenumi{\alph{enumi})}
\item
  If the codimension of Z is at least 2 at every point, the canonical injection of X into Y is relatively connected. Indeed, let z be a point of Z; there exists an open neighborhood V of z in Y and a homeomorphism \(\varphi\) of V onto a finite-dimensional vector space E over R such that \(\varphi(V \cap Z)\) is a vector subspace F of E whose codimension is \(\geqslant 2\) . The set E-F is connected (TG, VI, p. 5, proposition 4), whence the assertion.
\item
  Suppose moreover that the codimension of Z is at least 3 at every point; the hypotheses of Proposition 6 and its corollary are then satisfied, since, with the notation of the preceding paragraph, E-F is simply connected (I, p. 128, example 4).
\end{enumerate}

Since differentiable manifolds are loop-disentanglable spaces (IV, p. 347), the results of this section admit the following particular case.

Let Y be a locally finite-dimensional differentiable manifold, Z a closed submanifold of Y, and i the canonical injection of Y−Z into Y.

\begin{enumerate}
\def\labelenumi{\alph{enumi})}
\item
  If the codimension of Z in Y is ≥ 1 at every point, then the variety Y−Z is dense in Y and the map \(\pi_{0}(i)\) is surjective.
\item
  If the codimension of Z in Y is ≥ 2 at every point, the map \(\pi_{0}(i)\) is bijective and, for every point x in Y-Z, the map \(\pi_{1}(i,x)\) is surjective.
\item
  If the codimension of Z in Y is \(\geqslant 3\) at every point, the maps \(\pi_{0}(i)\) and \(\pi_{1}(i,x)\) are bijective, for every point x in Y-Z.
\end{enumerate}

\subsection{Presentation of the Poincaré groups}

THEOREM 1. --- Let X be a compact and loop-disentanglable topological space, and let x be a point of X. The Poincaré group \(\pi_{1}(X,x)\) is of finite presentation.

The arc-connected component of x in X is open, closed, and loop-disentanglable; this allows us to assume that the space X is connected. Since the space X is loop-disentanglable, the X-space \(E = \lambda_{x}(X)\) , equipped with the endpoint map, is a nonempty and simply path-connected covering (IV, p. 342, th. 1). The group \(G = \text{Aut}_{\text{X}}(\text{E})\) is isomorphic to \(\pi_{1}(\text{X}, x)\) ; it is therefore a matter of proving that the group G is finitely presented.

Since X is compact, every point x of X has a compact neighborhood \(K_{x}\) over which the covering E is trivializable (TG, I, p. 65, corollary). Since X is locally connected, every \(x \in X\) has an open neighborhood \(W_{x}\) , connected and contained in \(K_{x}\) . Let F be a finite subset of X such that the \(W_{x}\) , for \(x \in F\) , cover X. Let n be the cardinality of F.

Let us show by induction that there exists, for every integer \(k\) such that \(1 \leqslant k \leqslant n\) , a subset A of cardinality \(k\) contained in F and, for \(x \in A\) , a section \(s_x\) of \(p\) over \(K_x\) , such that the union of the \(s_x(W_x)\) , for \(x \in A\) , is a connected subset of E. The assertion is true for \(k = 1\) . Suppose it is true for an integer \(k\) such that \(1 \leqslant k < n\) and let us show that it is true for \(k + 1\) . Let A be a subset of F of cardinality \(k\) and, for all \(x \in A\) , let \(s_x\) be a section of E over \(K_x\) such that \(\bigcup_{x \in A} s_x(W_x)\) is connected. The open sets \(\bigcup_{x \in A} W_x\) and \(\bigcup_{x \in F-A} W_x\) are nonempty and cover X; their intersection is therefore nonempty, since X is connected. There thus exist \(x \in A\) , \(y \in F - A\) and \(z \in W_x \cap W_y\) . Set \(A' = A \cup \{y\}\) and let us choose a section \(s_y\) of E over \(K_y\) such that \(s_y(z) = s_x(z)\) . The connected open sets \(\bigcup_{p \in A} s_p(W_p)\) and \(s_y(W_y)\) are nonempty and have a point in common; their union is therefore connected. This proves the assertion for \(k + 1\) . By induction, it is therefore true for every integer \(k \in \{1, \ldots, n\}\) .

Applying the above to k = n , there then exists, for every \(x \in F\) , a section \(s_x\) of E over \(K_x\) , such that \(U = \bigcup_{x \in F} s_x(W_x)\) is an open connected subset of E. We have \(p(U) = X\) . Since the group G acts transitively in each fiber of the covering E, we have GU = E.

The closure of U is contained in \(\bigcup_{x\in\mathrm{F}}s_x(\mathrm{K}_x)\) , therefore is compact. The action of G on E is proper (I, p. 96, cor. 1), hence the set of pairs \((g,x)\in\mathrm{G}\times\mathrm{E}\) such that \(x\in\overline{\mathrm{U}}\) and \(gx\in\overline{\mathrm{U}}\) is a compact subset of \(\mathrm{G}\times\mathrm{E}\) (TG, I, p. 77, prop. 6). It follows that the set of \(g\in\mathrm{G}\) such that \(\overline{\mathrm{U}}\cap g\overline{\mathrm{U}}\neq\varnothing\) is compact. Since it is discrete, it is finite. A fortiori, the set of \(g\in\mathrm{G}\) such that \(\mathrm{U}\cap g\mathrm{U}\neq\varnothing\) is finite. The theorem thus follows from prop. 10 of I, p. 136.

\subsection{Complements on Polish spaces}

Lemma 1. --- Let X be a topological space and let A be a subset of X. For A to be meager, it is necessary and sufficient that there exists an open covering \((\mathrm{U}_{i})_{i\in\mathrm{I}}\) of X such that \(A\cap U_{i}\) be meager in \(U_{i}\) for all \(i\in I\) .

Let \(\mathcal{O}\) be the set of open subsets U of X such that \(A \cap U\) is meager. The family of subsets of \(\mathcal{O}\) consisting of pairwise disjoint open sets, ordered by inclusion, is inductive. Let \(\mathfrak{U}\) be a maximal element; such an element exists by E, III, p. 20, th. 2.

Let O be the union of the open sets of X belonging to \(\mathfrak{U}\). For every open set \(U \in \mathfrak{U}\) , let \((B_{U,n})\) be a sequence of nowhere dense subsets of U whose union is \(A \cap U\) . For every integer n, let \(B_n = \bigcup_{U \in \mathfrak{U}} B_{U,n}\) .

As U ranges over U, it is nowhere dense relative to O (TG, IX, p. 52, prop. 1). By TG, IX, p. 53, prop. 2, \(B_{n}\) is a nowhere dense subset of X, since O is open in X. Consequently, the set A ∩ O, equal to the union of the \(B_{n}\) is a meager subset of X.

Let F be the complement of O. It is a closed subset of X; let us show that it has empty interior. Otherwise, let x be an interior point of F. By hypothesis, there exists an open neighborhood V of x, which may be assumed to be contained in F, such that A ∩ V is meager. Then V is disjoint from the open sets of X belonging to U, and U ∪ \{V\} is a set of pairwise disjoint open sets belonging to O, which contradicts the maximality of U. The relation

\[
\mathrm{A} = (\mathrm{A} \cap \mathrm{F}) \cup (\mathrm{A} \cap \mathrm{O})
\]

then implies that A is meager, which is what was to be proved.

Let X be a topological space.

Recall (TG, IX, p. 69) that a subset A of X is said to be approachable if there exists an open set U of X such that \(U \cap \complement A\) and \(A \cap \complement U\) are meager in X. The collection of approachable subsets of X is a \(\sigma\)-algebra that contains the Borel \(\sigma\)-algebra (TG, IX, p. 69, Lemma 8 and its proof). A meager subset is approachable.

For any subset A of X, let D(A) be the set of points \(x \in X\) such that for every neighborhood V of x, A∩V is not meager. We also set \(\mathrm{D}^{*}(\mathrm{A}) = \mathrm{A} \cup \mathrm{D}(\mathrm{A})\) .

Lemma 2. --- Let X be a topological space and let A be a subset of X.

\begin{enumerate}
\def\labelenumi{\alph{enumi})}
\item
  The set D(A) is closed in X; its complement is the largest open set U in X such that A ∩ U is a meager set.
\item
  For A to be meager, it is necessary and sufficient that D(A) be empty.
\item
  The set \(\mathrm{D}^{*}(\mathrm{A})\) is approachable.
\item
  For every approachable set B of X containing A, \(D^{*}(A) \cap \complement B\) is meager.
\end{enumerate}

Let U be the union of the open sets V of X such that \(A \cap V\) is meager. For a point x to belong to U, it is necessary and sufficient that it have a neighborhood V such that \(A \cap V\) is meager. Thus \(U = \complement D(A)\) , and D(A) is therefore a closed subset of X. By construction, the subspace U has a covering by open sets whose intersection with A is meager. By Lemma 1, applied to the topological space U and the subset \(A \cap U\) , \(A \cap U\) is meager in U, hence also in X. This proves assertion a), since, by construction, every open set V of X such that \(A \cap V\) is meager is contained in U.

Assertion b) follows immediately from this.

\begin{enumerate}
\def\labelenumi{\alph{enumi})}
\setcounter{enumi}{2}
\item
  We have \(D^{*}(A) = A \cup D(A) = (A \cap U) \cup D(A)\) . The set \(D(A)\) is closed, hence approachable (IV, p. 361); the meager part \(A \cap U\) is also approachable. Consequently, \(D^{*}(A)\) is approachable (loc. cit.).
\item
  Let B be an approachable subset of X containing A. Its complement \(\complement B\) is then an approachable subset of X (loc. cit.), and there therefore exists an open set V of X such that \(V \cap B\) and \(\complement V \cap \complement B\) are meager. Since \(A \subset B\) , \(V \cap A\) is still meager, whence \(V \subset U\) . Since \(\complement U \cap \complement B\) is contained in \(\complement V \cap \complement B\) , it is also a meager subset of X. Finally, the inclusions
\end{enumerate}

\[
\mathrm{D} ^ {*} (\mathrm{A}) \cap \complement \mathrm{B} = ((\mathrm{A} \cap \mathrm{U}) \cap \complement \mathrm{B}) \cup (\complement \mathrm{U} \cap \complement \mathrm{B}) \subset (\mathrm{A} \cap \mathrm{U}) \cup (\complement \mathrm{U} \cap \complement \mathrm{B})
\]

prove that \(D^{*}(A) \cap \complement B\) is a meager subset of X.

Remark 1. --- Let G be a topological group and let A be a meager subset of G. Suppose that A contains a nonempty open set U of G, and let y be a point of U. Then U is meager, and for every \(x \in G\) , \(xy^{-1}U\) is a neighborhood of x in G. Consequently, every point of G has a meager neighborhood, so G is meager (IV, p. 360, lemma 1).

Conversely, if G is not meager, every meager subset has empty interior and G is a Baire space. Since a Baire space is not meager, this proves the following assertion: For a topological group to be a Baire space, it is necessary and sufficient that it be not meager.

PROPOSITION 7. --- Let X be a Hausdorff space. Every Souslin subspace (TG, IX, p. 59, def. 2) of X is approachable.

Let S be a Souslin subspace of X. By definition, there exists a complete metric space with countable base P and a continuous surjective map g from P onto S. By TG, IX, p. 64, lemma 3, there exists a sieve \(\mathrm{C} = (\mathrm{C}_{n}, p_{n}, \varphi_{n})_{n \in \mathbf{N}}\) of the metric space P.

For every integer \(n\) and every \(c \in \mathrm{C}_n\) , denote by \(\mathrm{F}_n(c)\) the image of the set \(\varphi_n(c)\) under \(g\) ; let us also set \(\mathrm{F}_n^*(c) = \mathrm{D}^*(\mathrm{F}_n(c))\) and

\[
\mathrm{G} _ {n} (c) = \mathrm{F} _ {n} ^ {*} (c) \cap \mathbb {C} \left(\bigcup_ {c ^ {\prime} \in p _ {n} ^ {- 1} (c)} \mathrm{F} _ {n + 1} ^ {*} (c ^ {\prime})\right).\tag{1}
\]

For every \(c \in C_{n}\) , \(\mathrm{F}_{n}^{*}(c)\) is approachable (IV, p. 361, lemma 2, c)). Since \(C_{n+1}\) is countable, the union of the approachable subsets \(\mathrm{F}_{n+1}^{*}(c')\) , for \(c'\) ranging over \(p_{n}^{-1}(c)\) , is again an approachable subset of X. It contains the union of the \(\mathrm{F}_{n+1}(c')\) , which equals \(\mathrm{F}_{n}(c)\) . Consequently (loc. cit., d)), \(\mathrm{G}_{n}(c)\) is a meager subset of X. The union G of the sets \(\mathrm{G}_{n}(c)\) , for \(n \in \mathbf{N}\) and \(c \in C_{n}\) , is therefore a meager subset of X.

Let \(c_0 \in \mathrm{C}_0\) and let \(x\) be an element of \(\mathrm{F}_0^*(c_0) \cap \complement \mathrm{G}\) ; let us prove that \(x \in \mathrm{F}_0(c_0)\) . Since \(x \notin \mathrm{G}_0(c_0)\) and \(x \in \mathrm{F}_0^*(c_0)\) , there exists \(c_1 \in p_0^{-1}(c_0)\) such that \(x \in \mathrm{F}_1^*(c_1)\) by relation (1). By induction, there exists an element \(c = (c_n)_n\) of \(\prod_n C_n\) such that, for every integer \(n \in \mathbf{N}\) , \(x \in \mathrm{F}_n^*(c_n)\) and \(p_n(c_{n+1}) = c_n\) .

The subsets \(\varphi_n(c_n)\) form a base of a Cauchy filter in P; this filter converges to a point \(p\) of P, since P is complete. The image of this filter under \(g\) is a filter F on X, with base the set of the \(\mathrm{F}_n(c_n)\) , for \(n\in \mathbf{N}\) , which converges to \(g(p)\) . Since \(\mathrm{F}_n^* (c_n)\) is contained in \(\overline{\mathrm{F}_n(c_n)}\) and \(x\) belongs to each of the \(\mathrm{F}_n(c_n)\) , \(x\) is an adherent point of F, hence \(x = g(p)\) since the space X is Hausdorff (TG, I, p. 52, prop. 1). The point \(p\) belongs to \(\overline{\varphi_1(c_1)}\) , hence to \(\varphi_0(c_0)\) . Consequently, \(x\in \mathrm{F}_0(c_0)\) .

Consequently, the set \(\mathrm{F}_{0}^{*}(c_{0}) - \mathrm{F}_{0}(c_{0})\) is contained in G. It is thus a meager subset of X, hence also an approachable subset. Since \(\mathrm{F}_{0}(c_{0}) = \mathrm{F}_{0}^{*}(c_{0}) - (\mathrm{F}_{0}^{*}(c_{0}) - \mathrm{F}_{0}(c_{0}))\) and \(\mathrm{F}_{0}^{*}(c_{0})\) is approachable, \(\mathrm{F}_{0}(c_{0})\) is approachable, because the approachable subsets of X form a \(\sigma\)-algebra. Since \(C_{0}\) is countable, the set S which is the union of the \(\mathrm{F}_{0}(c_{0})\) , for \(c_{0} \in C_{0}\) , is still approachable.

\subsection{Meager equivalence relations in Polish spaces}

Lemma 3. --- Let \((\mathrm{U}_{i})_{i\in\mathrm{I}}\) be a finite family of nonempty open sets of a topological space X and let O be an open dense subset of \(X\times X\) . There exists a family \((\mathrm{V}_{i})_{i\in\mathrm{I}}\) of nonempty open subsets of X such that \(V_{i}\subset U_{i}\) for all \(i\in I\) and such that \(V_{i}\times V_{i'}\subset O\) for every pair \((i,i')\) of distinct elements of I.

For every pair \((j_{1}, j_{2})\) of distinct elements of I, the set of families \((x_{i}) \in \mathrm{X}^{\mathrm{I}}\) such that \((x_{j_{1}}, x_{j_{2}}) \in \mathrm{O}\) is a dense open subset of \(X^{I}\) . The intersection \(\Omega\) of these open sets, as \((j_{1}, j_{2})\) ranges over the finite set of pairs of distinct elements of I, is therefore a dense open set of \(X^{I}\) .

Therefore, \(\Omega \cap \prod_{i\in \mathbf{I}}\mathrm{U}_i\) is a nonempty open set of \(\mathrm{X}^{\mathrm{I}}\) , hence contains an open set of \(\mathrm{X}^{\mathrm{I}}\) of the form \(\prod_{i\in \mathbf{I}}\mathrm{V}_i\) , where for every \(i\in \mathbf{I}\) , \(\mathrm{V}_i\) is a nonempty open subset of \(\mathrm{U}_i\) . This proves the lemma.

PROPOSITION 8. --- Let P be a nonempty Polish topological space and let R be an equivalence relation on P whose graph is a meager subset of P × P. There exists a continuous injective map from the set \(\{0,1\}^{\mathbf{N}}\) , endowed with the product topology of the discrete topologies, into P whose image meets each equivalence class under R in at most one point.

Endow the space P with a metric \(d\) compatible with its topology for which it is complete. Let \((\mathrm{A}_n)_{n\in \mathbf{N}}\) be an increasing sequence of closed subsets of \(\mathrm{P}\times \mathrm{P}\) , with empty interiors, such that the graph \(\Gamma_{\mathrm{R}}\) of the relation R is contained in the union of the \(\mathrm{A}_n\) .

Let \(\mathcal{O}\) be the set of nonempty open subsets of P. We shall construct by induction a sequence \((f_n)_{n \in \mathbf{N}}\) , where for every \(n\) , \(f_n\) is a map from \(\{0,1\}^n\) into \(\mathcal{O}\) , satisfying the following properties:

\begin{enumerate}
\def\labelenumi{(\roman{enumi})}
\item
  For every \(n \geqslant 1\) , every \(x \in \{0,1\}^{n-1}\) and every \(t \in \{0,1\}\) , the closure of the open set \(f_n(x,t)\) is contained in \(f_{n-1}(x)\) ;
\item
  For every \(x \in \{0,1\}^n\) , we have \(\mathrm{diam}(f_n(x)) \leqslant 2^{-n}\) ;
\item
  For every \(n \geqslant 1\) and every pair \((x, x')\) of distinct elements of \(\{0, 1\}^n\) , \(f_n(x) \times f_n(x')\) does not meet \(\mathrm{A}_{n-1}\) .
\end{enumerate}

Choose a nonempty open set U of P of diameter \(\leqslant 1\) and one defines \(f_{0}\) as the constant map with image \(\{U\}\) . Suppose the maps \(f_{0}, \ldots, f_{n}\) have been constructed.

Let \(p \colon \{0,1\}^{n+1} \to \{0,1\}^n\) be the map defined by \(p(x_0, \ldots, x_n) = (x_0, \ldots, x_{n-1})\) . By Lemma 3 above, applied to the family of open sets \((f_n(p(x)))\) for \(x \in \{0,1\}^{n+1}\) and to the dense open set \(\mathbb{C}\mathrm{A}_n\) of \(\mathrm{P} \times \mathrm{P}\) , there exists a family \((g(x))_{x \in \{0,1\}^{n+1}}\) of nonempty open subsets of P such that \(g(x) \subset f_n(p(x))\) for all \(x\) and such that \((g(x) \times g(x')) \cap \mathrm{A}_n\) is empty for every pair \((x,x')\) of distinct elements of \(\{0,1\}^{n+1}\) . One then defines the map \(f_{n+1}\) by choosing, for each element \(x \in \{0,1\}^{n+1}\) , a nonempty open subset of \(g(x)\) whose diameter is \(\leqslant 2^{-n-1}\) and whose closure is contained in \(g(x)\) .

For every element \(x = (x_n)_{n \in \mathbf{N}}\) of \(\{0,1\}^{\mathbf{N}}\) , the sequence of sets \((f_n(x_0, \ldots, x_{n-1}))_{n \in \mathbf{N}}\) is a decreasing sequence of open subsets of P each of which contains the closure of the next and whose diameter tends to 0; the intersection of this sequence of sets is therefore reduced to a point (TG, II, p. 15), which one denotes by \(f(x)\) . If two points \(x\) , \(x'\) of \(\{0,1\}^{\mathbf{N}}\) satisfy \(x_i = x'_i\) for \(i \leqslant n\) , we have \(d(f(x), f(x')) \leqslant 2^{-n}\) . Consequently, the map \(f: \{0,1\}^{\mathbf{N}} \to \mathbf{P}\) is continuous. Let \(x\) and \(x'\) be distinct elements of \(\{0,1\}^{\mathbf{N}}\) . For \(n \in \mathbf{N}\) such that \((x_0, \ldots, x_n) \neq (x_0', \ldots, x_n')\) , the open set \(f_{n+1}(x_0, \ldots, x_n) \times f_{n+1}(x_0', \ldots, x_n')\) is disjoint from \(A_n\) , by definition of \(f_{n+1}\) , so the pair \((f(x), f(x'))\) does not belong to \(A_n\) . It follows that \(f(x)\) and \(f(x')\) are not equivalent for the relation R. Consequently, \(f\) is injective and its image meets each equivalence class under R in at most one point. The proposition is thus proved.

\subsection{Cardinality of Poincaré groups}

PROPOSITION 9. --- Let X be a loop-disentanglable topological space and let W be a base of the topology of X. For every point x of X, the cardinality of the group \(\pi_{1}(X,x)\) is bounded above by sup(Card(W),Card(N)). In particular, the Poincaré group of a metric space of countable type that is loop-disentanglable is countable.

Indeed, the space \(\lambda_{x}(\mathrm{X})\) equipped with the endpoint map is a connected covering of X whose fiber at x is \(\pi_{1}(\mathrm{X}, x)\) . The assertion then follows from I, p. 40, th. 3.

Lemma 4. --- Let X be a Baire space, let G be a topological group acting continuously on X, and let B be an approachable subset of X that is not meager. The set of points \(g \in G\) such that \(B \cap gB \neq \emptyset\) is a neighborhood of the identity element of G.

Since B is approachable, \(\complement B\) is also approachable, because the collection of approachable subsets of X is a \(\sigma\)-algebra. There therefore exists an open subset U of X such that \(U \cap \complement B\) and \(B \cap \complement U\) are meager in X.

Let V be a neighborhood of the identity element of G and W a nonempty open subset of X contained in U such that \(V \cdot W \subset U\) . For all \(g \in V\) , \(U \cap gU\) is not empty.

Let \(g \in G\) such that \(B \cap gB\) is empty. The relations

\[
\begin{array}{l} \mathrm{U} \cap g \mathrm{U} = (\mathrm{U} \cap g \mathrm{U}) \cap \mathbb {C} (\mathrm{B} \cap g \mathrm{B}) \\ \qquad = (\mathrm{U} \cap g \mathrm{U}) \cap (\mathbb {C} \mathrm{B} \cup g \mathbb {C} \mathrm{B}) \\ \qquad = (\mathrm{U} \cap g \mathrm{U} \cap \mathbb {C} \mathrm{B}) \cup (\mathrm{U} \cap g \mathrm{U} \cap g \mathbb {C} \mathrm{B}) \\ \qquad \subset (\mathrm{U} \cap \mathbb {C} \mathrm{B}) \cup g (\mathrm{U} \cap \mathbb {C} \mathrm{B}) \end{array}
\]

imply that U ∩ gU is meager in X. Since it is an open subset and X is a Baire space, it is empty. We therefore have \(g \notin V\) , whence the lemma.

THEOREM 2 (Shelah \(^{(2)}\) ). --- Let X be a connected and locally path-connected Polish space and let x be a point of X. If X is not loop-disentanglable, the group \(\pi_{1}(X,x)\) has the cardinality of the continuum.

Let d be a metric defining the topology of X for which it is complete. Suppose that X is not loop-disentanglable. Then there exists a point \(a \in X\) and, for every integer \(n \geqslant 0\) , a loop \(c_{n}\) in X whose image has diameter \(\leqslant 2^{-n}\) and whose class in \(\pi_{1}(X, a)\) is not trivial.

Let K denote the set \(\{0,1\}^{\mathbf{N}}\) and endow it with the product topology, the space \(\{0,1\}\) equipped with the discrete topology. For every element \(\varepsilon = (\varepsilon_n)\) of K, let \(c_{\varepsilon}\) be the map from I to X defined by \(c_{\varepsilon}(0) = a\) and, for \(2^{-n - 1} \leqslant t \leqslant 2^{-n}\) , \(c_{\varepsilon}(t) = c_n(2^{n + 1}t - 1)\) if \(\varepsilon_n = 1\) and \(c_{\varepsilon}(t) = a\) otherwise. We have \(c_{\varepsilon}(0) = c_{\varepsilon}(1) = a\) . The map \(c_{\varepsilon}\) is continuous at every point \(t \in ]0,1]\) . It is also so at 0 because \(d(a,c_{\varepsilon}(t)) \leqslant 2^{-n}\) if \(t \in [0,2^{-n}] \) . The map \(c_{\varepsilon}\) is therefore a loop in X at \(a\) .

If \(\varepsilon\) and \(\varepsilon'\) are elements of K such that \((\varepsilon_0, \ldots, \varepsilon_n) = (\varepsilon_0', \ldots, \varepsilon_n')\) , then \(c_{\varepsilon}(t) = c_{\varepsilon'}(t)\) for all \(t \in [2^{-n-1}, 1]\) and \(d(c_{\varepsilon}(t), c_{\varepsilon'}(t)) \leqslant d(a, c_{\varepsilon}(t)) + d(a, c_{\varepsilon'}(t)) \leqslant 2^{-n}\) if \(t \in [0, 2^{-n-1}]\) . It follows that the map \(\varepsilon \mapsto c_{\varepsilon}\) from K into the space \(\Omega_a(X)\) is continuous, when the space \(\Omega_a(X)\) is equipped with the topology of compact convergence.

Denote by \(\Gamma \subset \mathrm{K} \times \mathrm{K}\) the set of pairs \((\varepsilon, \varepsilon')\) such that \(c_{\varepsilon}\) is strictly homotopic to \(c_{\varepsilon'}\) . This is the graph of an equivalence relation R on K.

Lemma 5. --- The set \(\Gamma\) is a meager subset of K × K.

Let Z denote the space \(K \times K \times \mathcal{C}_{c}(\mathbf{I} \times \mathbf{I}; X)\) . The topology of the space \(\mathcal{C}_{c}(\mathbf{I} \times \mathbf{I}; X)\) is defined by the metric \(\delta\) given by \(\delta(h, h') = \sup_{u \in \mathbf{I} \times \mathbf{I}} d(h(u), h'(u))\) and it is complete for this metric (TG, X, p. 20, corollary and TG, X, p. 9, cor. 1); by TG, X, p. 24, th. 1, this topology is of countable type. The space \(\mathcal{C}_{c}(\mathbf{I} \times \mathbf{I}; X)\) is therefore a Polish space. The same holds for the space Z, since K is a Polish space (TG, IX, p. 57, prop. 1).

Let H be the subset of Z consisting of the triples \((\varepsilon,\varepsilon^{\prime},h)\) such that h is a strict homotopy connecting \(c_{\varepsilon}\) to \(c_{\varepsilon^{\prime}}\) . The maps from Z into \(X^{2}\) given by \(a_{t}\colon(\varepsilon,\varepsilon^{\prime},h)\mapsto(c_{\varepsilon}(t),h(t,0))\) , \(b_{t}\colon(\varepsilon,\varepsilon^{\prime},h)\mapsto(c_{\varepsilon^{\prime}}(t),h(t,1))\) and \(c_{t}\colon(\varepsilon,\varepsilon^{\prime},h)\mapsto(h(0,t),h(1,t))\) are continuous, for every \(t\in I\) , because the map \(\varepsilon\mapsto c_{\varepsilon}\) from K to \(\Omega_{a}(X)\) is continuous, as are the mappings \(h\mapsto h(s,t)\) of \(\mathcal{C}_{c}(\mathbf{I}\times\mathbf{I};\mathbf{X})\) into X. By definition, H is the intersection of the sets \(a_{t}^{-1}(\Delta_{\mathrm{X}})\) , \(b_{t}^{-1}(\Delta_{\mathrm{X}})\) and \(c_{t}^{-1}((a,a))\) , for \(t\in I\) , where \(\Delta_{X}\) denotes the diagonal of X. This proves that H is a closed subset of Z.

Consequently, H is a Polish space. Let \(p \colon \mathrm{Z} \to \mathrm{K} \times \mathrm{K}\) be the canonical projection; we have \(\Gamma = p(\mathrm{H})\) by definition. Since \(\mathrm{K} \times \mathrm{K}\) is separated, \(\Gamma\) is a Souslin subset of \(\mathrm{K} \times \mathrm{K}\) (TG, IX, p. 59, def. 2). According to IV, p. 362, prop. 7, this is an approachable subset of \(\mathrm{K} \times \mathrm{K}\) .

Suppose that \(\Gamma\) is not meager. The space K × K is a compact topological space, hence a Baire space (TG, IX, p. 55, th. 1). Endow the group \(G_{0}\) of permutations of the set \(\{0,1\}\) with the discrete topology, and let the product topological group \(G = G_{0}^{\mathbf{N}}\) act diagonally on \(K = \{0,1\}^{\mathbf{N}}\) . The group G then acts continuously on K × K via the map (g, (x,y)) ↦ (x, g · y). By Lemma 4, the set V of elements \(g \in G\) such that \(\Gamma \cap g \cdot \Gamma \neq \varnothing\) is a neighborhood of the identity element of G.

Let \(g \in V\) ; let \((\varepsilon, \varepsilon') \in \Gamma \cap g \cdot \Gamma\) . Since \((\varepsilon, \varepsilon') \in \Gamma\) , we have \(\varepsilon\,R\,\varepsilon'\) ; since \((\varepsilon, \varepsilon') \in g \cdot \Gamma\) , \(g^{-1} \cdot (\varepsilon, \varepsilon') = (\varepsilon, g^{-1}\cdot\varepsilon') \in \Gamma\) , so that \(\varepsilon\,R\,g^{-1}\cdot\varepsilon'\) . We have thus shown that, for every \(g \in V\) , there exists \(\varepsilon \in K\) such that \(\varepsilon\) and \(g\varepsilon\) are equivalent for R.

For \(m \in \mathbf{N}\) , denote by \(\tau_m\) the element of G all of whose terms are equal to \(e\) except the one with index \(m\) which is equal to the nontrivial element \(\tau\) of \(G_0\) . There exists an integer \(m\) such that \(\tau_m\) belongs to V; then choose \(\varepsilon \in K\) such that \(\varepsilon\) and \(\varepsilon' = \tau_m \cdot \varepsilon\) are equivalent for R. This implies that the loops \(c_{\varepsilon}\) and \(c_{\varepsilon'}\) are strictly homotopic. By construction, these loops coincide on the intervals \([0, 2^{-m-1}]\) and \([2^{-m}, 1]\) ; on the interval \([2^{-m-1}, 2^{-m}]\) , one is the constant map with image a and the other is the map \(t \mapsto c_{m}(2^{m+1}t - 1)\) . It follows that \(c_{m}\) is strictly homotopic to the constant loop with image \(\{a\}\) , whence a contradiction. Lemma 5 is thus proved.

We now finish the proof of Theorem 2. By Prop. 8, there exists a continuous injective map \(\gamma\) of \(\{0,1\}^{\mathbf{N}}\) into K whose image meets every equivalence class under R in at most one point. If \(k\) and \(k'\) are distinct elements of \(\{0,1\}^{\mathbf{N}}\) , the loops \(c_{\gamma(k)}\) and \(c_{\gamma(k')}\) are not strictly homotopic in X and the map \(\{0,1\}^{\mathbf{N}} \to \pi_1(\mathrm{X},a)\) given by \(k \mapsto [c_{\gamma(k)}]\) is injective. In particular, \(\operatorname{Card}(\pi_1(\mathrm{X},a)) \geqslant \operatorname{Card}(\{0,1\}^{\mathbf{N}}) = \operatorname{Card}(\mathfrak{P}(\mathbf{N}))\) . Since X is a metrizable topological space of countable type, so is \(\Omega_a(\mathrm{X})\) (TG, X, p. 24, th. 1). Consequently, \(\Omega_a(\mathrm{X})\) is homeomorphic to a subspace of \([0,1]^{\mathbf{N}}\) (TG, IX, p. 18, prop. 12) and

\[
\begin{array}{r l} & {\mathrm{Card} (\Omega_ {a} (\mathbf {X})) \leqslant \mathrm{Card} ([ 0, 1 ] ^ {\mathbf {N}}) = \mathrm{Card} (\mathfrak {P} (\mathbf {N}) ^ {\mathbf {N}})} \\ & {\qquad = \mathrm{Card} (\mathfrak {P} (\mathbf {N} \times \mathbf {N})) = \mathrm{Card} (\mathfrak {P} (\mathbf {N})).} \end{array}
\]

A fortiori, \(\operatorname{Card}(\pi_{1}(\mathrm{X},a))\leqslant\operatorname{Card}(\mathfrak{P}(\mathbf{N}))\) . It then follows from cor. 2 of E, III, p. 25, that \(\pi_{1}(\mathrm{X},a)\) has the cardinality of the continuum, as was to be proved.

Example. --- Let P be the topological space which is the union of the circles with center \((2/n,0)\) passing through the origin of the plane \(R^{2}\) , for \(n \geqslant 1\) (III, p. 336, exerc. 6). The Poincaré group of P has the cardinality of the continuum (III, p. 338, exerc. 9).

\section{POINCARÉ GROUPS OF TOPOLOGICAL GROUPS}

\subsection{Extension of local homomorphisms of groups}

DEFINITION 1. --- Given a topological group G and a group G', a local homomorphism of G into G' is a map f from a neighborhood V of the identity element of G into G' such that, for every pair of points x, y in V with xy ∈ V, one has f(xy) = f(x)f(y).

If G' is a topological group and if f is a continuous map, one says that f is a continuous local homomorphism (or local morphism of topological groups).

If G and G' are topological groups, the notion of local isomorphism from G to G' was defined (TG, III, p. 6, def. 2). A local isomorphism from G to G' is a homeomorphism f from a neighborhood V of the identity element of G onto a neighborhood V' of the identity element of G' such that f and the inverse map of f are local homomorphisms.

PROPOSITION 1. --- Let G and G' be topological groups and let p: G → G' be a group homomorphism. For p to make G a covering of G', it is necessary and sufficient that the restriction of p to a suitable neighborhood of the identity element of G be a local isomorphism from G to G'.

The condition is necessary by proposition 3 of TG, III, p. 6. Conversely, suppose that \(p\) induces a homeomorphism from an open neighborhood V of the identity element \(e\) of G onto a neighborhood V' of the identity element of G'. The map \(p\) is then continuous (TG, III, p. 15, prop. 23) and open (TG, III, p. 16, prop. 24). The image H of \(p\) is a subgroup of G′ containing V, hence open and closed in G′ (TG, III, p. 7, corollary). Let N be the kernel of \(p\) ; we have N∩V = \{e\} ; hence N is discrete (loc. cit., prop. 5). Consequently, \(p\) makes G a covering of H, principal with group N (I, p. 100, cor. 3 of th. 1). Since H is open and closed in G', the G'-space (G, p) is a covering.

PROPOSITION 2. --- Let G be a connected topological group, let G' be a group, and let \(f: V \to G'\) be a local homomorphism of G into G', where V is a connected neighborhood of the identity element of G. Then there exists a connected topological group H, a morphism of topological groups \(p\colon\mathrm{H}\to\mathrm{G}\) such that \((\mathrm{H},p)\) is a covering of G and a group homomorphism \(\varphi\colon\mathrm{H}\to\mathrm{G}'\) such that the set of \(y\in p^{-1}(\mathrm{V})\) satisfying \(f(p(y))=\varphi(y)\) is a neighborhood of the identity element in H.

If G' is a topological group and the map f is continuous, such a homomorphism \(\varphi\) is continuous.

Lemma 1. — Let G be a topological group and let V be a connected neighborhood of the identity element e of G. For every neighborhood U of e in G and every \(x \in V\) , there exists an integer \(n \in N\) and elements \(u_{1}, \ldots, u_{n}\) in U such that \(u_{1} \ldots u_{n} = x\) and \(u_{1} \ldots u_{k} \in V\) for every integer k such that \(1 \leqslant k \leqslant n\) .

We may assume that U is open and contained in V. Let A denote the set of \(x \in V\) satisfying the condition of the lemma. If \(x \in A\) and if \(y \in xU \cap V\) then \(y \in A\) , whence \(AU \cap V \subset A\) ; this shows that A is open in V. Let \(x \in V\) such that \(xU^{-1} \cap A \neq \varnothing\) ; we then have \(x \in AU \cap V\) , therefore \(x \in A\) . Consequently, if \(x \in V\) and \(x \notin A\) , \(xU^{-1} \cap A = \varnothing\) and A is closed in V. Since \(e \in A\) and V is connected, it follows that A = V.

Let us now prove the proposition. Denote by \(j\) the map \(g \mapsto (g, f(g))\) from V to \(\mathbf{G} \times \mathbf{G}'\) and let H be the subgroup of \(\mathbf{G} \times \mathbf{G}'\) generated by \(j(\mathrm{V})\) .

Let U be a neighborhood of \(e\) in G, contained in V. Let \(x \in V\) ; by Lemma 1, there exists \(u_1, \ldots, u_n \in U\) such that \(x = u_1 \ldots u_n\) and such that \(u_1 \ldots u_k \in V\) for every integer \(k\) such that \(1 \leqslant k \leqslant n\) . By induction, we have \(f(u_1 \ldots u_k) = f(u_1) \ldots f(u_k)\) for every integer \(k\) , \(1 \leqslant k \leqslant n\) . In particular, \(j(x)\) belongs to the subgroup generated by \(j(U)\) . It follows that H is generated by \(j(U)\) .

Let \(\mathcal{B}\) be the set of subsets of H of the form \(j(\mathrm{U})\) , where U is a neighborhood of \(e\) in V. Let us show that there exists a unique topology on H, compatible with its group structure, for which \(\mathcal{B}\) is a base of the neighborhood filter of the identity element. For this, it suffices to show that the set \(\mathcal{B}\) satisfies the conditions \((\mathrm{GV}_{\mathrm{I}}^{\prime})\) , \((\mathrm{GV}_{\mathrm{II}}^{\prime})\) and \((\mathrm{GV}_{\mathrm{III}}^{\prime})\) from III, p. 4.

Let U be a neighborhood of e in G, contained in V; there exists a neighborhood \(U'\) of e such that \(U' \cdot U' \subset U\) . For every pair x, y of points of \(U'\) , we have \(xy \in V\) and \(f(xy) = f(x)f(y)\) . Consequently, \(j(\mathrm{U}') \in \mathcal{B}\) and \(j(\mathrm{U}') \cdot j(\mathrm{U}') \subset j(\mathrm{U})\) . This shows that the condition \((\mathrm{GV}_{\mathrm{I}}')\) is satisfied.

The set \(U'' = V \cap U^{-1}\) is then a neighborhood of e in V and, for \(x \in U''\) , we have \(x^{-1} \in U\) and \(f(x^{-1}) = f(x)^{-1}\) . Consequently, \(j(\mathrm{U}'')^{-1} \subset j(\mathrm{U})\) , which shows the condition \((\mathrm{GV}_{\mathrm{II}}' )\) .

Finally, let us fix a neighborhood U of e in V such that \(U = U^{-1}\) and \(U^{3} \subset V\) . Let W be a neighborhood of e in U and let \(h = (g, f(g))\) be an element of \(j(\mathrm{U})\) . There exists a neighborhood \(W'\) of e contained in W such that \(gW'g^{-1} \subset W\) . Then, \(j(\mathrm{W}') \in \mathcal{B}\) and \(hj(\mathrm{W}')h^{-1} \subset j(\mathrm{W})\) , since we have \(f(gxg^{-1}) = f(g)f(x)f(g^{-1})\) , for \(x \in W'\) .

Let \(h \in H\) . Since U is a neighborhood of e contained in V, \(j(\mathrm{U})\) generates H; since U is symmetric, there exist elements \(u_{1}, \ldots, u_{n}\) in U such that \(h = j(u_{1}) \ldots j(u_{n})\) . By induction on n, there exists a neighborhood \(W'\) of e contained in W such that \(hj(\mathrm{W}')h^{-1} \subset j(\mathrm{W})\) . Consequently, the condition \((\mathrm{GV}_{\mathrm{III}}')\) is satisfied.

We then endow the group H with this topology.

Denote by \(p \colon \mathrm{H} \to \mathrm{G}\) the restriction to H of the first projection \(\mathrm{G} \times \mathrm{G}' \to \mathrm{G}\) . This is a group homomorphism. For every neighborhood U of \(e\) contained in V, the set \(p^{-1}(\mathrm{U})\) contains the neighborhood \(j(\mathrm{U})\) of the identity element of H, hence \(p\) is a continuous homomorphism of topological groups. For every neighborhood U of \(e\) contained in V, we have \(p(j(\mathrm{U})) = \mathrm{U}\) , therefore \(p\) is an open mapping. Since G is connected, \(p\) is surjective. Its kernel is discrete because it does not meet \(j(\mathrm{V})\) except at the identity element. It then follows from Corollary 3 (I, p. 100) that the G-space \((\mathrm{H}, p)\) is a covering.

Let \(\varphi\) be the restriction to H of the second projection \(G \times G' \to G'\) . If \(g \in V\) , we have \((g, f(g)) = j(g) \in j(V)\) and \(\varphi(g, f(g)) = f(g)\) , so that \(\varphi(y) = f(p(y))\) for \(y \in j(V)\) .

Suppose further that G' is a topological group and that the map f is continuous at e. Then the homomorphism \(\varphi\) is continuous at the identity element of H, hence is continuous (TG, III, p. 15, prop. 23).

COROLLARY 1. --- Let G be a simply connected topological group and let G' be a group. Let V be a connected neighborhood of the identity element in G and let f: V → G' be a local homomorphism. There exists a unique group homomorphism h: G → G' extending f. If G' is a topological group and if the map f is continuous, then the homomorphism h is continuous.

The group G is connected. Let H, \(\varphi\), and p be as in Proposition 2. Since G is simply connected, H is a trivializable covering of G (I, p. 124, def. 3); since H is connected and nonempty, the map \(p\) is an isomorphism of topological groups. Set \(h = \varphi \circ p^{-1}\) ; the map \(h\) is a group homomorphism and there exists an open neighborhood U of the identity element \(e\) of G contained in V such that \(f|U = h|U\) . It follows from Lemma 1 that \(f\) and \(h\) coincide on V. In other words, the map \(h\) extends the map \(f\) .

If G' is a topological group and the map f is continuous, the homomorphism \(\varphi\) is continuous, so h is too.

Let us prove the uniqueness of such an extension. The set of points of G where two homomorphisms from G to G′ coincide is a subgroup of G. Since the group G is connected, every subgroup of G containing a neighborhood of the identity element is equal to G (TG, III, p. 8, prop. 6). The uniqueness of h follows from this.

Remark 1. --- When G = R, this corollary follows from Proposition 6 of TG, V, p. 3.

COROLLARY 2. --- Two locally connected and simply connected topological groups that are locally isomorphic are isomorphic.

Let G and G' be locally connected topological groups, let V and V' be connected neighborhoods of the identity element of G and of G', respectively, and let \(f: V \to V'\) be a homeomorphism that is a local isomorphism from G to G'. By Corollary 1, there exists a unique continuous group homomorphism \(\varphi: G \to G'\) which extends f and a unique continuous group homomorphism \(\varphi': G' \to G\) which extends \(f^{-1}\) . The homomorphisms \(\varphi' \circ \varphi\) and \(\operatorname{Id}_{\mathrm{G}}\) coincide on a neighborhood of e in G, hence are equal, since G is connected. Likewise, \(\varphi \circ \varphi' = \operatorname{Id}_{\mathrm{G}'}\) , whence the corollary.

COROLLARY 3. --- Let G be a simply connected topological group and let V be a connected neighborhood of the identity element of G. One defines a presentation of G by taking as generating set the set V and as the set \(\mathbf{r}\) of relators the family of \(xyz^{-1}\) , where \((x,y,z)\) traverses the triples of elements of V such that \(xy = z\) .

Let \(\mathrm{F}(\mathrm{V}, \mathbf{r})\) be the quotient group of the free group \(\mathrm{F}(\mathrm{V})\) by the smallest normal subgroup containing the elements \(xyz^{-1}\) , where \((x, y, z) \in \mathrm{V} \times \mathrm{V} \times \mathrm{V}\) and xy = z (A, I, p. 86). Let us denote by \(f: V \to F(V, r)\) and \(g: F(V, r) \to G\) the canonical maps. By construction, the map f is a local homomorphism from G to \(\mathrm{F}(V, \mathbf{r})\) . By Corollary 1, there exists a unique group homomorphism \(\overline{f}: G \to F(V, r)\) extending f.

Since the group \(F(V, \mathbf{r})\) is generated by \(f(V)\) , the homomorphism \(\overline{f}\) is surjective. For \(x \in V\) , we have \(g(\overline{f}(x)) = g(f(x)) = x\) ; since V generates G, \(g \circ \overline{f} = \operatorname{Id}_{\operatorname{G}}\) , which proves that \(\overline{f}\) is injective. It is therefore an isomorphism.
